% Generated draft: requires mathematical and source review.
\documentclass[11pt]{article}
\usepackage[margin=25mm]{geometry}
\usepackage{fontspec}
\setmainfont{Latin Modern Roman}
\usepackage{amsmath,amssymb,mathtools,unicode-math}
\setmathfont{Latin Modern Math}
\usepackage{enumitem,needspace,xurl,hyperref}
\setlength{\parindent}{0pt}
\setlength{\parskip}{5pt}
\setlength{\emergencystretch}{3em}
\newcommand{\N}{\mathbb{N}}
\newcommand{\Z}{\mathbb{Z}}
\newcommand{\Q}{\mathbb{Q}}
\newcommand{\R}{\mathbb{R}}
\newcommand{\C}{\mathbb{C}}
\newcommand{\F}{\mathbb{F}}
\newcommand{\A}{\mathbb{A}}
\newcommand{\PP}{\mathbb{P}}
\newcommand{\E}{\mathbb{E}}
\newcommand{\eps}{\varepsilon}
\newcommand{\dd}{\,\mathrm{d}}
\begin{document}
\section*{UnsolvedMath ID: 30004520}
\textbf{Source problem code:} catalogue entry
\subsection{Definitions and mathematical statement}
In the usual Euclidean interpretation of the source, let \(d\geq 1\) be an integer and let \(S=\operatorname{conv}\{v_0,\ldots,v_d\}\subset\R^d\) be a nondegenerate simplex, meaning that its vertices are affinely independent. The source does not explicitly specify the ambient metric; the following construction uses Euclidean distance.

A longest-edge bisection of \(S\) chooses indices \(i\ne j\) such that
\[
\lVert v_i-v_j\rVert=\max_{0\leq a<b\leq d}\lVert v_a-v_b\rVert,
\]
sets \(m=(v_i+v_j)/2\), and replaces \(S\) by
\[
S_i=\operatorname{conv}\bigl(\{m\}\cup\{v_k:k\ne i\}\bigr),\qquad
S_j=\operatorname{conv}\bigl(\{m\}\cup\{v_k:k\ne j\}\bigr).
\]
These children have disjoint interiors, their union is \(S\), and each has half its \(d\)-dimensional volume. A descendant is a simplex obtained by finitely many such bisections. For a nondegenerate Euclidean simplex \(T\), let \(r(T)\) be its inradius and \(R(T)\) its circumradius, and define its shape ratio by
\[
q(T)=\frac{r(T)}{R(T)}.
\]

The supplied question asks whether, starting from a sufficiently generic \(d\)-dimensional simplex and repeatedly applying longest-edge bisection, the resulting subdivision stabilizes after finitely many iterations for every starting simplex in every dimension, or whether simplices produced from some fixed starting simplex can have shape ratios tending to zero.

Under an interpretation that recursively refines every child, the degeneration alternative can be expressed as the existence of a dimension \(d\), an admissible initial simplex \(S\), and descendants \(T_n\) of that same \(S\) such that
\[
\lim_{n\to\infty}\frac{r(T_n)}{R(T_n)}=0.
\]
Here admissibility remains subject to the source's unspecified genericity requirement. The source does not say whether the \(T_n\) must lie along one nested branch of descendants or may be selected from different branches.

A complete quantified stabilization assertion cannot be recovered from this record: ``sufficiently generic'' and ``stabilize'' are not defined, the refinement schedule and treatment of tied longest edges are unspecified, and the opening genericity restriction is not reconciled with the later phrase ``every starting simplex.'' In particular, stabilization must not silently be identified with finitely many similarity classes or with a positive lower bound for \(q(T)\). The supplied record also does not establish that failure of its intended stabilization property is equivalent to the stated degeneration alternative.
\subsection{Short English statement}
Repeatedly split a simplex at the midpoint of a longest edge, producing two equal-volume children. The source asks whether this refinement eventually stabilizes in every dimension, or whether descendants of one starting simplex can have inradius-to-circumradius ratio tending to \(0\). The meanings of stabilization, sufficient genericity, and the refinement schedule require clarification before the question has a fully precise formulation.
\subsection{Sources}
\begin{itemize}
\item Local UnsolvedMath registry: supplied complete record for permanent ID \(30004520\), titled ``Stability of Longest Edge Simplex Refinement,'' in the Oberwolfach Reports Open Problems collection; exact source problem code \texttt{catalogue entry}. The record lists proposed year 2020 and names no proposer.
\item Supplied report citation: \emph{Discrete Geometry} (2021), Oberwolfach Reports, DOI 10.4171/owr/2020/30, pp.\ 1469--1528. \url{https://doi.org/10.4171/owr/2020/30}. The supplied background describes the reports as documenting mathematical workshops and their open-problem discussions.
\item Supplied literature source: Rote, \emph{Open Problems in Discrete Geometry}. \url{https://page.mi.fu-berlin.de/rote/Kram/Problems-Discrete-Geometry-2020.pdf}. This attribution identifies a reference, not a named proposer of this problem.
\item Supplied literature source: \emph{Oberwolfach Reports 2020/30}. \url{https://ems.press/content/serial-article-files/46867}.
\item Literature assessment dated exactly \texttt{2026-08-22}: the supplied curation labels the problem open and states that stabilization is proved in dimensions two and three and for special four-dimensional starting simplices, while the question for all simplices in all dimensions remains open. This is supplied dated curation, not an independently verified literature assessment. The supplied background explicitly requests independent verification of time-sensitive status and cited claims.
\item The registry labels the statement \texttt{corrected\_verified}, explaining that only formula delimiters, notation, grammar, and typography were repaired and that mathematical content and scope were preserved. This supplied editorial label does not constitute mathematical verification of this draft.
\end{itemize}
\end{document}
